%======================= PAGES PRÉLIMINAIRES ===========================
\begin{titlepage}
\centering
\thispagestyle{empty}

% Bandeau institutionnel : l'établissement à gauche, la filière à droite.
% Ces deux logos priment visuellement sur celui du laboratoire, placé plus bas
% à côté de sa dénomination.
\noindent
\begin{minipage}[c]{0.46\textwidth}
  \raggedright
  \includegraphics[height=16mm]{\figpath/logos/fsbm.png}
\end{minipage}
\hfill
\begin{minipage}[c]{0.46\textwidth}
  \raggedleft
  \includegraphics[height=16mm]{\figpath/logos/dsbd.png}
\end{minipage}

\vspace{9mm}

{\large\scshape Université Hassan II de Casablanca\par}
{\large\scshape Faculté des Sciences Ben M'Sik\par}
\vspace{4mm}
{\normalsize Master Data Science \& Big Data\par}
{\normalsize Année universitaire 2025--2026\par}

\vspace{11mm}
{\color{liasrule}\rule{\textwidth}{0.8pt}}
\vspace{8mm}

{\Large\bfseries\color{liasblue} Projet de fin de module\par}
\vspace{3mm}
{\Large\bfseries\color{liasblue} Java Enterprise Edition\par}

\vspace{8mm}
{\Huge\bfseries\color{liasnavy} Plateforme de gestion\\[3mm] du laboratoire LIAS\par}

% Logotype officiel du laboratoire, adossé à sa dénomination.
\vspace{6mm}
\includegraphics[height=17mm]{\figpath/logos/lias.png}

\vspace{2mm}
{\large\itshape Laboratoire d'Intelligence Artificielle\\ et Systèmes\par}

\vspace{7mm}
{\color{liasrule}\rule{\textwidth}{0.8pt}}

\vspace{14mm}
{\normalsize Site institutionnel public et système d'information interne\\
pour la gestion des membres, de la gouvernance, de l'activité\\
scientifique et de la mémoire administrative du laboratoire\par}

\vfill

\begin{minipage}{0.46\textwidth}
\raggedright
{\small\bfseries\color{liasgrey} Réalisé par\par}
\vspace{1mm}
{\normalsize Douae Moussaoui\par}
\end{minipage}
\hfill
\begin{minipage}{0.46\textwidth}
\raggedleft
{\small\bfseries\color{liasgrey} Encadrant\par}
\vspace{1mm}
{\normalsize Pr.~Abdessamad Belangour\par}
\end{minipage}

\vspace{12mm}
{\small\color{liasgrey} Casablanca --- Septembre 2026\par}

\end{titlepage}

%----------------------------------------------------------------------
\cleardoublepage
\pagenumbering{roman}
\pagestyle{plain}

\chapter*{Remerciements}
\addcontentsline{toc}{chapter}{Remerciements}
\markboth{Remerciements}{}

Je tiens à exprimer ma reconnaissance à mon encadrant, \textbf{Pr.~Abdessamad Belangour},
pour la qualité de son encadrement tout au long du module Java Enterprise Edition et pour un
énoncé de projet suffisamment exigeant pour imposer de véritables décisions d'architecture
plutôt qu'un simple exercice de programmation.

Mes remerciements vont également au laboratoire LIAS, dont le fonctionnement réel ---
statuts d'appartenance, mandats de gouvernance, mémoire documentaire --- a servi de
cadre métier à ce travail et lui a donné sa consistance.

Je remercie enfin l'ensemble du corps enseignant du Master pour la formation dispensée
tout au long de l'année.

%----------------------------------------------------------------------
\chapter*{Résumé}
\addcontentsline{toc}{chapter}{Résumé}
\markboth{Résumé}{}

Ce rapport présente la conception et la réalisation d'une plateforme web de gestion pour
le laboratoire de recherche LIAS, développée dans le cadre du module Java Enterprise
Edition du Master Data Science \& Big Data. L'application remplit simultanément deux
fonctions : celle de site institutionnel public et celle de système d'information interne
couvrant l'ensemble du cycle de vie administratif et scientifique du laboratoire.

La difficulté centrale du sujet n'est pas la quantité de modules mais une exigence
structurelle du cahier des charges : la séparation stricte entre le \srt{} d'appartenance
d'un membre, son \man{} de gouvernance et le \rol{} applicatif qui gouverne ses
permissions. Cette contrainte a orienté toute la conception. Le rôle n'est jamais saisi
manuellement : il est \emph{dérivé} du statut et des mandats actifs par un service
dédié, historisé à chaque changement, et réconcilié automatiquement au démarrage de
l'application.

La solution repose sur une architecture trois tiers : un backend Spring~Boot~3.3
(Java~21) exposant 22 contrôleurs REST adossés à 30 entités JPA et à une base
PostgreSQL versionnée par Flyway ; une interface React~19 de 21 pages construite sur un
système de composants interne ; une couche de sécurité stateless fondée sur JWT,
BCrypt et un contrôle d'accès par rôle appliqué méthode par méthode. La traçabilité est
assurée par un journal d'audit et quatre tables d'historique, conformément à la règle du
cahier des charges selon laquelle aucune donnée n'est jamais perdue.

La validation combine 66 tests automatisés (aucun échec) et un audit fonctionnel de bout
en bout portant sur 21 routes et les 9 rôles du système. Les 21 sections du cahier des
charges sont couvertes ; les limites résiduelles, notamment l'absence de thème sombre et
de mesure de couverture de code, sont explicitement identifiées.

\vspace{4mm}
\noindent\textbf{Mots-clés} --- Java Enterprise Edition ; Spring Boot ; JPA / Hibernate ;
React ; PostgreSQL ; JWT ; contrôle d'accès par rôle (RBAC) ; gouvernance de laboratoire ;
traçabilité ; historisation ; architecture trois tiers.

%----------------------------------------------------------------------
\chapter*{Abstract}
\addcontentsline{toc}{chapter}{Abstract}
\markboth{Abstract}{}

This report describes the design and implementation of a web management platform for the
LIAS research laboratory, carried out as the final project of the Java Enterprise Edition
module of the Data Science \& Big Data Master's programme. The application serves two
purposes at once: a public institutional website and an internal information system
covering the whole administrative and scientific lifecycle of the laboratory.

The core difficulty of the assignment is not the number of modules but a structural
requirement of the specification: the strict separation between a member's membership
\emph{status}, their governance \emph{mandate}, and the application \emph{role} that
drives their permissions. This constraint shaped the entire design. The role is never set
by hand: it is \emph{derived} from status and active mandates by a dedicated service,
historised on every change, and automatically reconciled at application start-up.

The solution follows a three-tier architecture: a Spring~Boot~3.3 backend (Java~21)
exposing 22 REST controllers over 30 JPA entities and a PostgreSQL database versioned with
Flyway; a React~19 front end of 21 pages built on an in-house component system; and a
stateless security layer based on JWT, BCrypt and method-level role-based access control.
Traceability relies on an audit journal and four history tables, in line with the
specification rule that no data is ever lost.

Validation combines 66 automated tests (no failures) with an end-to-end functional audit
covering 21 routes and the 9 system roles. All 21 sections of the specification are
covered; remaining limitations --- notably the absence of a dark theme and of a measured
code-coverage figure --- are explicitly stated.

\vspace{4mm}
\noindent\textbf{Keywords} --- Java Enterprise Edition; Spring Boot; JPA / Hibernate;
React; PostgreSQL; JWT; role-based access control; laboratory governance; traceability;
auditability; three-tier architecture.

%----------------------------------------------------------------------
\cleardoublepage
\tableofcontents
\cleardoublepage
\listoffigures
\addcontentsline{toc}{chapter}{Liste des figures}
\cleardoublepage
\listoftables
\addcontentsline{toc}{chapter}{Liste des tableaux}

%----------------------------------------------------------------------
\cleardoublepage
\chapter*{Liste des acronymes}
\addcontentsline{toc}{chapter}{Liste des acronymes}
\markboth{Liste des acronymes}{}

\begin{longtable}{@{}p{2.6cm}p{11.6cm}@{}}
\toprule
\textbf{Acronyme} & \textbf{Signification} \\
\midrule
\endfirsthead
\toprule
\textbf{Acronyme} & \textbf{Signification} \\
\midrule
\endhead
API      & \emph{Application Programming Interface} --- interface de programmation \\
BCrypt   & Fonction de hachage de mots de passe à coût paramétrable \\
CdC      & Cahier des charges \\
CDP      & \emph{Chrome DevTools Protocol} \\
CORS     & \emph{Cross-Origin Resource Sharing} \\
CRUD     & \emph{Create, Read, Update, Delete} \\
DTO      & \emph{Data Transfer Object} --- objet de transfert \\
E2E      & \emph{End-to-End} --- de bout en bout \\
HS512    & HMAC-SHA-512, algorithme de signature des jetons JWT \\
JEE      & \emph{Java Enterprise Edition} (aujourd'hui Jakarta EE) \\
JPA      & \emph{Jakarta Persistence API} \\
JPQL     & \emph{Jakarta Persistence Query Language} \\
JSON     & \emph{JavaScript Object Notation} \\
JWT      & \emph{JSON Web Token} \\
LIAS     & Laboratoire d'Intelligence Artificielle et Systèmes \\
ORM      & \emph{Object-Relational Mapping} --- correspondance objet-relationnel \\
PV       & Procès-verbal (de réunion) \\
RBAC     & \emph{Role-Based Access Control} --- contrôle d'accès fondé sur les rôles \\
REST     & \emph{Representational State Transfer} \\
SGBD     & Système de gestion de base de données \\
SPA      & \emph{Single-Page Application} \\
SMTP     & \emph{Simple Mail Transfer Protocol} \\
SQL      & \emph{Structured Query Language} \\
UML      & \emph{Unified Modeling Language} \\
UUID     & \emph{Universally Unique Identifier} \\
\bottomrule
\end{longtable}

\cleardoublepage
\pagenumbering{arabic}
\setcounter{page}{1}
\pagestyle{fancy}
